\section{Methodology}
\subsection{Attention to distance}
For each layer $\ell$ we collect the head-averaged attention matrix
$G^{(\ell)}\in\mathbb{R}^{T\times T}$, symmetrise it by
$G \leftarrow \max(G, G^\top)$, and convert to a metric:
\begin{equation}
d^{(\ell)}_{ij} = 1 - G^{(\ell)}_{ij}, \qquad d^{(\ell)}_{ii}=0.
\end{equation}

\subsection{Vietoris--Rips filtration}
We feed the distance matrix $d^{(\ell)}$ to \texttt{ripser}, computing
the persistence diagram in dimensions $0$ and $1$. From each diagram
we extract:
\begin{itemize}
\item \textbf{Total persistence} in $H_k$: $\sum_i (\text{death}_i - \text{birth}_i)$
over all finite bars.
\item \textbf{Betti curve} $\beta_k(t)$: the number of bars alive at
each filtration threshold $t \in \{0, \tfrac{1}{32}, \dots, 1\}$.
\end{itemize}

\subsection{Statistical test}
For each layer we run a Welch two-sample $t$-test on
$\text{total persistence}(H_0)$ between the truthful and hallucinated
conditions. We report the per-layer $p$-value as
$-\log_{10} p$.
